\documentclass[10pt,twocolumn]{article}
\usepackage[T1]{fontenc}
\usepackage[utf8]{inputenc}
\usepackage{lmodern}
\usepackage[margin=1.8cm,top=2.4cm,bottom=2cm,columnsep=0.6cm]{geometry}
\usepackage{fancyhdr}
\usepackage{titlesec}
\usepackage{booktabs}
\usepackage{amsmath,amssymb,amsfonts}
\usepackage{microtype}
\usepackage{xcolor}
\usepackage{mdframed}
\usepackage{parskip}
\usepackage{enumitem}
\usepackage{hyperref}

\definecolor{rulered}{RGB}{180,30,30}
\definecolor{boxgray}{RGB}{245,245,245}
\definecolor{keywordgray}{RGB}{100,100,100}

\pagestyle{fancy}
\fancyhf{}
\fancyhead[L]{\textsc{\small Claude Mag}}
\fancyhead[C]{\small\textit{Machine Learning Research Digest}}
\fancyhead[R]{\small 2026-09-14}
\fancyfoot[C]{\small\thepage}
\renewcommand{\headrulewidth}{0.8pt}

\titleformat{\section}{\normalsize\bfseries\color{rulered}}{}{0pt}{}%
  [\vspace{-4pt}\color{rulered}\rule{\columnwidth}{0.4pt}]
\titlespacing{\section}{0pt}{8pt}{4pt}

\hypersetup{colorlinks=true,urlcolor=rulered,linkcolor=black}

\begin{document}
\setlength{\parindent}{0pt}
\setlength{\parskip}{4pt}

{\small\color{keywordgray}\textbf{Keywords:}
  distribution shift \textbullet{}
  covariate shift \textbullet{}
  concept shift \textbullet{}
  optimal transport \textbullet{}
  generalization bounds}\par\vspace{4pt}

{\LARGE\bfseries\raggedright
  General Quantification of Covariate
  and Concept Shifts}\par\vspace{4pt}

{\large\itshape\raggedright
  Entropic Optimal Transport Yields the First General,
  Estimable Error Bound Unifying Both Shift Types}\par\vspace{6pt}

{\small\textbf{By} Hongbo Chen, Li Charlie Xia
  (South China University of Technology)}\par\vspace{2pt}

{\small\color{keywordgray}arXiv:2609.11918 \textbullet{} ICML 2026}
\par\vspace{2pt}\noindent{\color{rulered}\rule{\columnwidth}{0.6pt}}\vspace{6pt}

Generalization under distribution shift is a core unsolved challenge in
modern machine learning. Existing learning bound theory handles covariate
shift and concept shift in separate, narrow, idealized frameworks ---
and both frameworks break when the source and target distributions do
not share the same support, which is the common case in deployment. This
paper derives a single general bound that unifies both shift types using
entropic optimal transport, is estimable from samples, and comes with an
algorithm --- DataShifts --- that measures each component of the shift
from finite data.

\begin{mdframed}[backgroundcolor=boxgray,linecolor=rulered,linewidth=1pt,
    innerleftmargin=6pt,innerrightmargin=6pt,
    innertopmargin=6pt,innerbottommargin=6pt]
\textbf{\small AT A GLANCE}\par\vspace{4pt}
\begin{description}[leftmargin=0pt,labelsep=4pt,itemsep=2pt,parsep=0pt,
    font=\small\bfseries]
  \item[Problem:] Existing distribution-shift bounds are restricted to
    matching supports and separate shift types; neither condition holds
    in practice.
  \item[Why it matters:] Unestimable, vacuous, or type-specific bounds
    cannot guide real deployment decisions under shift.
  \item[Method:] Define $\gamma^*$-concept shifts via entropic OT;
    derive a unified bound; construct sample estimators for both components.
  \item[Result:] DataShifts produces finite non-vacuous bounds on all
    12 synthetic and 5 real benchmarks; 1.5--5$\times$ tighter than
    prior methods where prior methods are finite.
  \item[Evidence:] Formal concentration guarantees; ablations on SCM
    size, sample size, and cost function.
\end{description}
\end{mdframed}

\section{The Core Problem}

A learning algorithm is trained on source data $(X_s, Y_s) \sim P_s$
and deployed on target data $(X_t, Y_t) \sim P_t$. How much does the
test error on $P_t$ exceed the training error on $P_s$?

\textbf{Covariate shift} assumes the label mechanism is unchanged
($P_t(Y|X) = P_s(Y|X)$) but the input distribution shifts
($P_t(X) \neq P_s(X)$). Importance-weighted ERM corrects this.

\textbf{Concept shift} assumes the conditional itself changes:
$P_t(Y|X) \neq P_s(Y|X)$. Existing bounds measure this change via
$f$-divergences between conditionals --- but these divergences are
infinite when $\mathrm{supp}(P_t) \not\subseteq \mathrm{supp}(P_s)$,
producing vacuous bounds in the mismatch setting that characterizes
most real deployments.

\section{Why Existing Bounds Fail}

The authors prove a formal impossibility: any bound of the form
\[
  \mathcal{E}_t(h) \leq \mathcal{E}_s(h) + D_f(P_s, P_t) + \lambda
\]
where $D_f$ is any $f$-divergence-based concept shift measure will be
infinite whenever $\mathrm{supp}(P_t) \not\subseteq \mathrm{supp}(P_s)$,
even if the conditional distributions are in practice very close.

This is not merely a theoretical inconvenience. On standard domain
adaptation benchmarks such as Camelyon17 (where hospital scanners differ)
or CIFAR-10-C (where corruptions shift the input distribution globally),
the target support extends beyond the source support. All prior bounds
are vacuous on these benchmarks.

\section{Entropic Optimal Transport and $\gamma^*$-Concept Shifts}

The key insight is to replace $f$-divergence with a soft coupling via
entropic optimal transport. Given a cost function $c(x, x')$ on the
input space, define:
\[
  \gamma^* = \arg\min_{\gamma \in \Pi(P_s^X, P_t^X)}
  \mathbb{E}_\gamma[c(X,X')] + \varepsilon H(\gamma)
\]
where $H(\gamma)$ is the entropy of the coupling and $\varepsilon > 0$
is a regularization parameter. The \textbf{$\gamma^*$-concept shift} is:
\[
  \Delta_\gamma = \mathbb{E}_{\gamma^*}\!\left[
    \mathrm{TV}\!\left(P_s(Y|X{=}x),\, P_t(Y|X'{=}x')\right)
  \right]
\]
This measures how much the label distribution changes between
``matched'' source and target points under the optimal transport plan.
When supports agree, this recovers the classical TV-based concept
shift. When supports mismatch, the transport plan couples each target
point to the nearest available source point, producing a
well-defined, finite quantity.

\section{The General Error Bound}

\textbf{Theorem.} For any hypothesis $h$ with bounded loss
$\ell \in [0,1]$:
\[
  \mathcal{E}_t(h) \leq \mathcal{E}_s(h)
    + \underbrace{W_\varepsilon(P_s^X, P_t^X)}_{\text{covariate shift}}
    + \underbrace{\Delta_\gamma(P_s, P_t)}_{\gamma^*\text{-concept shift}}
    + O\!\left(\sqrt{\frac{\log|\mathcal{H}|}{n}}\right)
\]
where $W_\varepsilon$ is the entropic Wasserstein distance between input
marginals. This bound:
\begin{itemize}[itemsep=2pt,topsep=2pt]
  \item Applies to multiclass, regression, and stochastic labeling.
  \item Is finite for arbitrary source/target supports.
  \item Recovers known tight bounds as special cases.
  \item Is estimable from finite samples.
\end{itemize}

\section{The DataShifts Algorithm}

DataShifts estimates both shift components from source and target samples:

\textbf{Step 1.} Solve $\hat{\gamma}^* = \mathrm{Sinkhorn}(\hat{P}_s^X,
\hat{P}_t^X, c, \varepsilon)$ via Sinkhorn iterations.

\textbf{Step 2.} Compute $\hat{W}_\varepsilon$ as the primal OT objective.

\textbf{Step 3.} For each matched pair $(x_i, x_j')$ in $\hat{\gamma}^*$,
estimate $\mathrm{TV}(P_s(Y|x_i), P_t(Y|x_j'))$ from label empiricals.

\textbf{Concentration guarantee:} With probability $1 - \delta$,
\[
  |\hat{W}_\varepsilon - W_\varepsilon| \leq C\sqrt{\frac{\log(1/\delta)}{n}}
\]
and similarly for $\hat{\Delta}_\gamma$, where $C$ depends on the
diameter of the input space.

\section{Experimental Findings}

Experiments span 12 synthetic shift scenarios and 5 real-world
benchmarks (Portraits, CIFAR-10-C, Camelyon17, DomainNet, and a
tabular clinical dataset):

\begin{center}
\begin{tabular}{lcc}
\toprule
Dataset & $\mathcal{H}\Delta\mathcal{H}$ bound & DataShifts \\
\midrule
CIFAR-10-C (blur)  & $\infty$ & 0.48 (true: 0.31) \\
Camelyon17         & 2.31     & 0.39 (true: 0.24) \\
Portraits          & 0.87     & 0.27 (true: 0.18) \\
\bottomrule
\end{tabular}
\end{center}

DataShifts is 1.5--5$\times$ tighter where prior methods are finite,
and produces finite bounds where prior methods give $\infty$.

\section{Shift Decomposition}

A key practical capability of DataShifts is decomposing total shift
into covariate and concept components, guiding which correction strategy
to apply:

\textbf{CIFAR-10-C (corruption-only):}
$\hat{W}_\varepsilon = 0.41$, $\hat{\Delta}_\gamma = 0.03$ ---
shift is almost entirely covariate; apply importance weighting.

\textbf{Clinical tabular (changed deployment policy):}
$\hat{W}_\varepsilon = 0.07$, $\hat{\Delta}_\gamma = 0.28$ ---
shift is almost entirely concept; apply label calibration.

\section{Ablations and Interpretation}

\textbf{Regularization $\varepsilon$:} Smaller values tighten bounds
but increase optimization difficulty; $\varepsilon = 0.1$ works well
across all tested datasets.

\textbf{Sample size:} DataShifts requires roughly $n = 500$ target
samples for reliable bounds; $n = 100$ suffices for decomposition
quality.

\textbf{Cost function:} Euclidean cost works for image benchmarks;
learned feature-space cost improves results on tabular data.

\section*{Reference}

Hongbo Chen, Li Charlie Xia.
\textbf{General Quantification of Covariate and Concept Shifts.}
arXiv:2609.11918, ICML 2026.
\url{https://arxiv.org/abs/2609.11918}

\end{document}
